% Options for packages loaded elsewhere
\PassOptionsToPackage{unicode}{hyperref}
\PassOptionsToPackage{hyphens}{url}
%
\documentclass[
  11pt,
]{article}
\usepackage{amsmath,amssymb}
\usepackage{iftex}
\ifPDFTeX
  \usepackage[T1]{fontenc}
  \usepackage[utf8]{inputenc}
  \usepackage{textcomp} % provide euro and other symbols
\else % if luatex or xetex
  \usepackage{unicode-math} % this also loads fontspec
  \defaultfontfeatures{Scale=MatchLowercase}
  \defaultfontfeatures[\rmfamily]{Ligatures=TeX,Scale=1}
\fi
\usepackage{lmodern}
\ifPDFTeX\else
  % xetex/luatex font selection
\fi
% Use upquote if available, for straight quotes in verbatim environments
\IfFileExists{upquote.sty}{\usepackage{upquote}}{}
\IfFileExists{microtype.sty}{% use microtype if available
  \usepackage[]{microtype}
  \UseMicrotypeSet[protrusion]{basicmath} % disable protrusion for tt fonts
}{}
\makeatletter
\@ifundefined{KOMAClassName}{% if non-KOMA class
  \IfFileExists{parskip.sty}{%
    \usepackage{parskip}
  }{% else
    \setlength{\parindent}{0pt}
    \setlength{\parskip}{6pt plus 2pt minus 1pt}}
}{% if KOMA class
  \KOMAoptions{parskip=half}}
\makeatother
\usepackage{xcolor}
\usepackage[margin=1in]{geometry}
\usepackage{longtable,booktabs,array}
\usepackage{calc} % for calculating minipage widths
% Correct order of tables after \paragraph or \subparagraph
\usepackage{etoolbox}
\makeatletter
\patchcmd\longtable{\par}{\if@noskipsec\mbox{}\fi\par}{}{}
\makeatother
% Allow footnotes in longtable head/foot
\IfFileExists{footnotehyper.sty}{\usepackage{footnotehyper}}{\usepackage{footnote}}
\makesavenoteenv{longtable}
\setlength{\emergencystretch}{3em} % prevent overfull lines
\providecommand{\tightlist}{%
  \setlength{\itemsep}{0pt}\setlength{\parskip}{0pt}}
\setcounter{secnumdepth}{-\maxdimen} % remove section numbering
\ifLuaTeX
  \usepackage{selnolig}  % disable illegal ligatures
\fi
\usepackage{bookmark}
\IfFileExists{xurl.sty}{\usepackage{xurl}}{} % add URL line breaks if available
\urlstyle{same}
\hypersetup{
  pdftitle={Boundary Excess Calculus v2.0},
  pdfauthor={Htet Ko Ko Naing, Independent Researcher, ORCID: 0009-0000-6140-0495},
  hidelinks,
  pdfcreator={LaTeX via pandoc}}

\title{Boundary Excess Calculus v2.0}
\usepackage{etoolbox}
\makeatletter
\providecommand{\subtitle}[1]{% add subtitle to \maketitle
  \apptocmd{\@title}{\par {\large #1 \par}}{}{}
}
\makeatother
\subtitle{Final Monograph Assembly - Verified Kernel Release M1-M15}
\author{Htet Ko Ko Naing, Independent Researcher, ORCID:
0009-0000-6140-0495}
\date{2026-07-02}

\begin{document}
\maketitle

\section{Formal Claim Boundary}\label{formal-claim-boundary}

Boundary Excess Calculus v2.0 (BEC v2.0) is a boundary-threshold
calculus with a staged Lean 4 / Mathlib verification program. The
current release is a verified kernel chain, not a claim that every
theorem in the full monograph has been mechanized.

The formal claim of this release is:

\begin{quote}
BEC v2.0 has a Lean 4.31.0 / Mathlib checked milestone chain M1-M14,
with a final M15 publication and reproducibility package containing pass
logs, blocked-token scans, SHA256 ledgers, and release documentation.
\end{quote}

The following stronger claims are explicitly not made:

\begin{enumerate}
\def\labelenumi{\arabic{enumi}.}
\tightlist
\item
  BEC does not replace ZFC, HoTT, topology, measure theory, probability
  theory, stochastic analysis, or category theory.
\item
  BEC is not presented as an unrestricted universal solver.
\item
  The entire monograph is not claimed to be fully mechanized.
\item
  External solver outputs are trusted only after passing typed
  certificate-import witnesses.
\item
  Stochastic certificates are not silently promoted to deterministic
  certificates.
\end{enumerate}

\section{1. Foundational Framework and Universe
Core}\label{foundational-framework-and-universe-core}

A BEC v2.0 universe is an eleven-tuple

\[
B=(X, Ref, V,\le_V,0_V,\top_V,op,E,Assump,Claim,Cert).
\]

Here \(X\) is the carrier of objects, \(Ref\) is the type of reference
boundaries, and \(V\) is the value object into which boundary excess is
measured. The relation \(\le_V\) is a preorder on \(V\), \(0_V\) is the
zero-excess element, \(\top_V\) is a maximal or failure-tolerant bound,
and \(op:V\times V\to V\) is the additive or compositional excess
operation used by licensed aggregation. The excess map is

\[
E:Ref\times X\to V.
\]

The last three components are logical: \(Assump\) is the type of
assumptions, \(Claim\) is the type of BEC-native claims, and
\(Cert:Claim\to Type\) is the certificate family. A certificate is not a
proof of an arbitrary proposition; it is a typed witness for a
BEC-native claim under explicit routing discipline.

For \(A\in Ref\) and \(r,s\in V\), the basic threshold and shell
operators are

\[
T_r(A)=\{x\in X:E(A,x)\le_V r\},
\]

\[
T^{<}_r(A)=\{x\in X:E(A,x)<_V r\},
\]

\[
S_{r,s}(A)=\{x\in X:r<_V E(A,x)\le_V s\}.
\]

A metric instance takes \(V=[0,\infty]\) and

\[
E(A,x)=\inf_{a\in A}d(x,a).
\]

A topological Boolean instance takes \(V=\{0,1\}\), where \(0\) means
closure membership and \(1\) records excess from closure membership.

\section{2. Boundary-Threshold Graph
Layer}\label{boundary-threshold-graph-layer}

A Boundary-Threshold Graph is a five-tuple

\[
G_B=(N,E_G,\ell,\gamma,\omega).
\]

The node type \(N\) consists of typed claims. The edge relation
\(E_G\subseteq N\times N\) records claim-routing pathways. The map
\(\ell\) records the target claim level, \(\gamma\) gives the assumption
gate required to traverse an edge, and \(\omega\) is the obstruction
record emitted when the gate is not licensed.

A route is safe only when each edge is licensed. Thus BEC route
transport has the form

\[
\operatorname{LicensedRoute}(G_B,\Gamma,p,q)
\Rightarrow Cert(p)\to Cert(q).
\]

If a license is missing, BEC does not promote the claim. It reports an
obstruction. This is the central discipline that persists from v1.0
through v2.0.

\section{3. Computational Engine}\label{computational-engine}

BEC v2.0 restricts computational claims to BEC-shaped problems. A
problem is BEC-shaped when it is expressible as a boundary, threshold,
shell, profile, or certificate-routing claim over explicit carriers,
references, metrics, orders, or finite profiles.

\subsection{\texorpdfstring{3.1 Polyhedral \(L_\infty\)
case}{3.1 Polyhedral L\_\textbackslash infty case}}\label{polyhedral-l_infty-case}

Let

\[
A=\{y\in\mathbb{Q}^n:My\le b\}.
\]

For \(x\in\mathbb{Q}^n\), the claim

\[
E_\infty(A,x)\le r
\]

is represented by the linear feasibility system

\[
\exists y\in\mathbb{Q}^n\;\bigl(My\le b\;\wedge\;\forall i,\; -r\le x_i-y_i\le r\bigr).
\]

Equivalently, using slack variables \(u_i\ge 0\), it is enough to
require

\[
x_i-y_i\le u_i,
\qquad
 y_i-x_i\le u_i,
\qquad
u_i\le r.
\]

Then

\[
|x_i-y_i|\le r
\]

for every coordinate \(i\).

\subsection{\texorpdfstring{3.2 Polyhedral \(L_1\)
case}{3.2 Polyhedral L\_1 case}}\label{polyhedral-l_1-case}

The \(L_1\) threshold is encoded by slack variables

\[
u_i\ge 0,
\qquad x_i-y_i\le u_i,
\qquad y_i-x_i\le u_i,
\qquad \sum_i u_i\le r.
\]

This implies

\[
\sum_i |x_i-y_i|\le r.
\]

M8 verifies these two certificate shapes in Lean over rational vectors.
M13 extends the layer with a backend-import interface: a solver export
becomes a BEC certificate only when a Lean-side import witness supplies
status acceptance, norm agreement, radius agreement, polyhedron
membership, nonnegative slack variables, and the appropriate slack
inequalities.

\section{4. Compiler Normal Forms}\label{compiler-normal-forms}

The compiler layer is not a universal mathematical compiler. It is a
typed front end for the fragment of inputs that already lie within BEC's
boundary-threshold vocabulary. Its soundness statement is:

\[
\operatorname{normalizeProblem}(P)=\operatorname{ok}(N)
\Rightarrow
\operatorname{NativeNormalForm}(N).
\]

Thus successful compilation produces only BEC-native normal forms.
Failed compilation is not a defect; it is part of the safety boundary.
The compiler refuses to erase assumptions, silently coerce stochastic
claims into deterministic ones, or create unlicensed certificate
transport.

\section{5. Stochastic Boundary
Calculus}\label{stochastic-boundary-calculus}

Let \((\Omega,\mathcal F,\mathbb P)\) be a probability space. A
stochastic boundary claim has the form

\[
\mathbb P(E(A,Y)\le \tau)\ge \alpha.
\]

For two events

\[
P=\{E(A,Y)\le \tau\},\qquad Q=\{E(Y,B)\le \sigma\},
\]

if \(P\cap Q\subseteq R\), then the union bound gives

\[
\mathbb P(R)
\ge \mathbb P(P\cap Q)
=1-\mathbb P(P^c\cup Q^c)
\ge 1-\mathbb P(P^c)-\mathbb P(Q^c)
\ge \alpha+\beta-1.
\]

This is the verified M3 stochastic-transfer kernel. Empirical
certificates follow the Hoeffding shape

\[
\mathbb P(|\widehat p-p|\le \varepsilon)\ge 1-2\exp(-2n\varepsilon^2),
\]

and the deterministic Lean-side certificate rule records only the
algebraic consequence

\[
|\widehat p-p|\le\varepsilon,
\qquad \alpha\le \widehat p-\varepsilon
\Rightarrow
\alpha\le p.
\]

The exponential confidence expression is analytic data; the verified
core checks the implication once the deviation witness is supplied.

\section{6. Dynamic Boundary Layer}\label{dynamic-boundary-layer}

The dynamic layer records generator inequalities of the form

\[
\partial_t\phi+\mathcal L_t\phi\le -\kappa\phi+c.
\]

A dynamic risk certificate records a positive radius \(r>0\), moment
bound \(M\ge 0\), and safety parameter \(\alpha\) satisfying

\[
\alpha\le 1-\frac{M}{r}.
\]

M9 verifies that the moment-radius safety inequality is recovered from
the certificate data and that cut-locus local-time corrections are
explicit nonnegative bookkeeping terms. The release does not claim a
full mechanization of Itô-Tanaka analysis; it verifies the structural
certificate layer required to prevent unlicensed dynamic boundary
claims.

\section{7. Boolean Topology Fragment}\label{boolean-topology-fragment}

For a Kuratowski closure operator \(C\), define Boolean excess by

\[
E_C(A,x)=0 \iff x\in C(A),
\qquad
E_C(A,x)=1 \iff x\notin C(A).
\]

The zero threshold then recovers closure:

\[
T_0(A)=\{x:E_C(A,x)\le 0\}=C(A).
\]

M4 verifies this Boolean closure-excess fragment. The claim is
conservative: topology is represented as a Boolean BEC fragment; BEC is
not asserted to replace topology or its foundations.

\section{8. Boundary Type Theory}\label{boundary-type-theory}

Boundary Type Theory (BTT) is the type-theoretic layer for BEC-native
judgments. Its primitive forms include context formation, boundary type
formation, reference formation, excess judgment, threshold judgment,
shell judgment, route judgment, certificate judgment, and obstruction
judgment.

A typical soundness rule is

\[
Cert_\Gamma(\varphi)\Rightarrow \Gamma\models\varphi.
\]

M5 verifies the minimum BTT carrier, primitive judgment type,
certificate soundness, obstruction honesty construction, and
conservative interpretation soundness. BTT is a layer for threshold
reasoning; it is not a replacement for all type theory.

\section{9. Categorical Comparison
Layer}\label{categorical-comparison-layer}

A claim kernel consists of claims, certificates, routes, route identity,
route composition, and certificate transport. For routes

\[
\rho:p\to q,
\qquad
\sigma:q\to r,
\]

BEC requires

\[
transport(\sigma\circ\rho,c)=transport(\sigma,transport(\rho,c)).
\]

A kernel morphism preserves certificate transport:

\[
F(transport_K(\rho,c))=transport_L(F\rho,Fc).
\]

M14 verifies the categorical morphism, lax comparison, conservative
pullback, graph paths, and graph-path certificate transport. This is the
precise categorical comparison claim licensed by the release.

\section{10. Lean Verification Ledger}\label{lean-verification-ledger}

The recorded environment is

\[
\text{Lean }4.31.0,\qquad \text{Mathlib/Lake},\qquad \text{aarch64-unknown-linux-gnu}.
\]

Each checked milestone package contains a Lean file, shell check script,
pass log, blocked-token scan, and SHA256 evidence. The blocked-token
scan checks for

\[
\texttt{sorry},\quad \texttt{admit},\quad \texttt{axiom},\quad \texttt{unsafe}.
\]

The complete release ledger is included in the M15 publication package.

\section{11. Milestone Status}\label{milestone-status}

\begin{longtable}[]{@{}lll@{}}
\toprule\noalign{}
Milestone & Scope & Status \\
\midrule\noalign{}
\endhead
\bottomrule\noalign{}
\endlastfoot
M1 & Universe + Graph + Licensed Route Kernel & PASS \\
M2 & Compiler Normal Form Soundness & PASS \\
M3 & Probability / Stochastic Certificate Kernel & PASS \\
M4 & Topology Boolean Closure Kernel & PASS \\
M5 & Boundary Type Theory Foundation Kernel & PASS \\
M6 & Audit Manifest M1-M5 & GENERATED \\
M7 & All-in-One Kernel M1-M5 & PASS \\
M8 & Polyhedral Solver Kernel & PASS \\
M9 & Dynamic Generator Kernel & PASS \\
M10 & Final Audit Bundle M1-M9 & GENERATED \\
M11 & Monograph Integration Package & GENERATED \\
M12 & Extended All-in-One Kernel M1-M5 + M8 + M9 & PASS \\
M13 & Solver Backend Importer & PASS \\
M14 & Categorical Morphism / Lax Comparison Kernel & PASS \\
M15 & Publication / Reproducibility Package & GENERATED \\
\end{longtable}

\section{12. Package SHA256 Ledger}\label{package-sha256-ledger}

The full package SHA256 ledger is stored in the accompanying M15 ledger
and M16 machine-readable JSON file. The compact table below gives SHA256
prefixes for visual audit in the PDF edition.

\begin{longtable}[]{@{}
  >{\raggedright\arraybackslash}p{(\columnwidth - 4\tabcolsep) * \real{0.3333}}
  >{\raggedright\arraybackslash}p{(\columnwidth - 4\tabcolsep) * \real{0.3333}}
  >{\raggedright\arraybackslash}p{(\columnwidth - 4\tabcolsep) * \real{0.3333}}@{}}
\toprule\noalign{}
\begin{minipage}[b]{\linewidth}\raggedright
Milestone
\end{minipage} & \begin{minipage}[b]{\linewidth}\raggedright
Artifact
\end{minipage} & \begin{minipage}[b]{\linewidth}\raggedright
SHA256 prefix
\end{minipage} \\
\midrule\noalign{}
\endhead
\bottomrule\noalign{}
\endlastfoot
M1 & Universe + Graph + Licensed Route Kernel &
\texttt{52719cb28ec77997} \\
M2 & Compiler Normal Form Soundness & \texttt{d0c230856cbd0a05} \\
M3 & Probability / Stochastic Certificate Kernel &
\texttt{e1a76a74d612483a} \\
M4 & Topology Boolean Closure Kernel & \texttt{8a68bc95eee2eccc} \\
M5 & Boundary Type Theory Foundation Kernel &
\texttt{a497c7668fafa52d} \\
M6 & Audit Manifest M1-M5 & \texttt{660bde6fe57e4498} \\
M7 & All-in-One Kernel M1-M5 & \texttt{f58f7c861b1f78ea} \\
M8 & Polyhedral Solver Kernel & \texttt{231180aaecaf372d} \\
M9 & Dynamic Generator Kernel & \texttt{f70e426f01c4eacf} \\
M10 & Final Audit Bundle M1-M9 & \texttt{6cee643452a78a80} \\
M11 & Monograph Integration Package & \texttt{d5df73b30f93cd9d} \\
M12 & Extended All-in-One Kernel & \texttt{486bae681a3f2392} \\
M13 & Solver Backend Importer & \texttt{69a21a366e3d20f9} \\
M14 & Categorical Morphism Kernel & \texttt{20dc9851cad6d1a9} \\
M15 & Publication Package & \texttt{306c3fe0e2f59486} \\
\end{longtable}

\section{13. Release Statement}\label{release-statement}

BEC v2.0 is now assembled as a Lean-supported, audit-backed
boundary-threshold framework with verified milestone kernels for
universe formation, graph-based route licensing, compiler normal forms,
stochastic certificate transfer, Boolean topology, Boundary Type Theory,
polyhedral slack certificates, dynamic generator certificates, backend
certificate importing, and categorical comparison.

The release is intentionally conservative. Its strength is the explicit
distinction between drafted mathematical architecture, checked Lean
kernels, external solver evidence, stochastic analytic assumptions, and
publication claims.

\end{document}
